\subsection{An explicit harmonic twist with a nonzero global obstruction}

\begin{statement}{A.134}{Proved}
{A locally soluble harmonic quartic and completed old coordinates obstruct one puncture}
Put
\[
 \begin{gathered}
 h=\frac12,\qquad R=\frac{529}{382},\qquad
 c=-\frac{338}{191},\qquad q=139,\qquad r=1119169,\\
 Q=qr=155564491,\qquad
 a_0=\frac{529}{382},\quad a_1=\frac{147}{382},\quad
 a_h=\frac{(13Q)^2}{191}
       =\frac{4089852535355054689}{191},\\
 C:\quad T=a_0X_0^4,\quad T-1=a_1X_1^4,\quad
 T-h=a_hX_h^4,\qquad X_0X_1X_h\ne0,\\
 \rho=\frac{34}{7Q},\qquad
 d=7\cdot17Q=18512174429.
 \end{gathered}
\tag{A.134.1}
\]
The curve $C$ has ordinary points with $T\ne c$
over every completion of $\Q$.
Its smooth proper compactification has genus $33$.
The actual common field at the new boundary $T=c$ is
\[
 B=\Q(i,\sqrt6,\sqrt{299},\sqrt d),\qquad
 K=\Q(i,\sqrt6),\qquad [B:\Q]=16.
\tag{A.134.2}
\]
The cyclic class
\[
                \alpha=(T-c,d)\in\operatorname{Br}(\overline C)
\tag{A.134.3}
\]
has invariant $1/2$ at two and zero at every other
place, at every local point of $\overline C$.
Consequently $\overline C(\Q)$ is empty.
The canonical obstruction $\gamma_{C,c}$ of A.128
fixes $i,\sqrt6,\sqrt{299}$ and negates $\sqrt d$.
There is no $b\in\Q^\times$ for which the actual
augmentation $T-c=bY^2$ over $C$, with $Y\ne0$,
has points at all places.

There also exist completed classes
$A_0,A_1,A_h\in\widehat{\Q^\times}$ and ordinary
coordinates
\[
 t_p\in\Q_p\setminus\{0,1,h,c\}\quad(p<\infty),
            \qquad t_2=31/2,
\tag{A.134.4}
\]
with the exact localization identities
\[
 (A_0)_p=\kappa_p(t_p),\qquad
 (A_1)_p=\kappa_p(t_p-1),\qquad
 (A_h)_p=\kappa_p(t_p-h),\qquad
                    A_i\bmod4=[a_i].
\tag{A.134.5}
\]
All three signs are positive, with real centre $R$.
Every $A_i$ has infinitely many nonzero ordinary
integer prime valuations.
There is no rational squareclass matching
$[t_p-c]$ at all finite primes, and therefore no
completed new coordinate matching these same values.
Thus the nonzero global constant of A.128 occurs
for actual completed old data on the harmonic source.
\end{statement}
\begin{proof}
\emph{The exact cover, common character, and connected proper model.}
The two numbers $139$ and $1119169$ are prime;
the certificate below verifies primality by trial
division through the square root.
At $T=c$ the three radicands are
\[
 \begin{aligned}
 D_0&=-\frac{676}{529}
           =-\left(\frac{26}{23}\right)^2,\\
 D_1&=-\frac{1058}{147}
           =-6\left(\frac{23}{21}\right)^2,\\
 D_h&=-\frac{867}{2(13Q)^2}
           =-6\left(\frac{17}{26Q}\right)^2,\qquad
                  D_0D_1D_h=-\rho^2.
 \end{aligned}
\tag{A.134.6}
\]
The individual squareclasses generate $K=\Q(i,\sqrt6)$.
Among the seven nonempty subsets the only negative
rational-square products are the singleton $D_0$
and the full triple.
Their additional A.96 characters are respectively
$52/23$ and $2\rho$.
The literal square comparisons are
\[
             \frac{52/23}{299}=\left(\frac2{23}\right)^2,
          \qquad \frac{2\rho}{d}=\left(\frac2{7Q}\right)^2.
\tag{A.134.7}
\]
A.96 therefore proves exactly (A.134.2).
The four squareclasses $-1,6,299,d$ are independent:
their sign and prime supports distinguish them
($299=13\cdot23$ and $d=7\cdot17\cdot139\cdot r$).
This also proves the displayed degree.

Every factor contains the character field by an
explicit root expression.
In the actual fibre algebra put $z=Z_0Z_1Z_h$.
Then
\[
       z^4=-\rho^2,\qquad
       \iota=z^2/\rho,\qquad \iota^2=-1,\qquad
       \left(\frac{7Q}{2}(1-\iota)z\right)^2=d.
\tag{A.134.8}
\]
This is defined in each factor and proves the
common-character assertion without selecting one.
The corresponding singleton expression, scaled
by $23/2$, gives a square root of $299$.

Over $\overline\Q(T)$ the three radicands
$(T-i)/a_i$ are independent modulo fourth powers:
their valuations at the three different finite
punctures detect the three exponents modulo four.
The geometric degree is therefore $64$.
Each of $0,1,h,\infty$ has inertia order four and
sixteen geometric points above it.
For infinity this follows, for example, from the
local coordinates $z=1/X_0$, $Y_i=X_i/X_0$:
\[
       T=a_0z^{-4},\qquad
                 Y_i^4=(a_0-i z^4)/a_i\quad(i=1,h).
\tag{A.134.9}
\]
Thus Riemann--Hurwitz gives
$2g-2=-128+4\cdot64(1-1/4)=64$, or $g=33$,
and there are $64$ old geometric boundary points.

For an odd prime away from the coefficient and
denominator support, the three finite punctures
have distinct reductions and all coefficients are
units.
At a finite boundary $i$, $X_i$ is a local
parameter and the other equations read
$X_j^4=(i-j+a_iX_i^4)/a_j$.
Their constant terms are units and their root
derivatives are units.
The same assertion at infinity is (A.134.9).
On the remaining open part it is the derivative
$4X_i^3$.
These charts give a smooth proper tame reduction
with the same degree, inertia and genus.
This specifies the good-reduction model used below.

\emph{Actual local points, including every good prime.}
The complete support set for the coefficients,
new-boundary differences and character is
\[
                         S=\{2,3,7,13,17,23,139,191,r\}.
\tag{A.134.10}
\]
Useful rational centres at these places are
\[
 \begin{array}{c|c}
 p&T_p\\ \hline
 2&31/2\\
 3,7,13,23,191&R\\
 17&1+a_1\,17^4=12277969/382\\
 139&h+(169/191)139^2=6530689/382\\
 r&h+a_h=8179705070710109569/382.
 \end{array}
\tag{A.134.11}
\]
The retained data file gives one rational $T_p$
at every prime $p\le4500$, together with the last
row at $r$.
There are $611$ centres in total, and each avoids
$0,1,h,c$.
For each of the three ratios the independent
certificate verifies the exact fourth-power criterion
\[
 \begin{cases}
 v_2((T_p-i)/a_i)\equiv0\pmod4,\quad
       \text{normalized unit }=1\pmod{16},&p=2,\\
 v_p((T_p-i)/a_i)\equiv0\pmod4,\quad
       \bar u^{(p-1)/\gcd(4,p-1)}=1,&p\text{ odd}.
 \end{cases}
\tag{A.134.12}
\]
These are sufficient as well as necessary:
at odd $p$ the residue fourth root lifts because
its derivative is a unit; at two the unit
fourth-power map is exactly $1+16\Z_2$, as proved
in (A.127.7)--(8).
Multiplying these unit roots by the indicated
integral powers of $p$ gives actual roots in $\Q_p$.
Thus the certificate supplies local points, not
merely congruences of approximate coordinates.

For every remaining good prime, Hasse--Weil gives
at least $p+1-66\sqrt p$ points on the proper
reduction; the exact curve-cohomology and trace
inputs are
\cite[Proposition 14.2]{AgentMilneLEC} and
\cite[(1.5.1), Theorem (1.6)]{AgentDeligneWeilI}.
For $p\ge4500$ one has
\[
                         p+1-66\sqrt p-64>0.
\tag{A.134.13}
\]
Indeed $(4500-63)^2-66^2\cdot4500=84969>0$,
and the squared difference is strictly increasing
thereafter.
Removing the $64$ old boundary points therefore
leaves an ordinary residue point with all three
root coordinates nonzero.
The unit derivatives lift it.
If its residue centre is $c$, choose a different
lift of that centre in the same residue class and
lift the three roots again; this gives $T_p\ne c$.
Only the old boundary was removed in (A.134.13);
the unramified fibre over $c$ is handled by this
explicit change of lift.

At the real place take $T=R$.
The first two fourth-power ratios are one, and
the third is $Q^{-2}>0$.
This proves ordinary local solubility at all places.

\emph{The exact global invariant and the finite obstruction.}
A.128 extends $\alpha$ to the proper curve and
proves that its evaluation is constant at every
place.
We compute it using (A.134.11).
At two,
\[
          d\equiv5\pmod8,\qquad
          T_2-c=\frac{6597}{382},\qquad v_2(T_2-c)=-1.
\tag{A.134.14}
\]
The field $\Q_2(\sqrt d)=\Q_2(\sqrt5)$ is
unramified quadratic, and its norm group has
even valuations.
Hence the invariant is $1/2$.

At $3,13,23$ both the first entry and $d$ are
units, so the invariant is zero.
At $191$ the character is locally split: the
exact residue check gives $(d/191)=1$.
At the three ramified character primes $7,17,139$,
the first-entry residues are respectively
\[
        (T_7-c)\bmod7=2=3^2,\quad
        (T_{17}-c)\bmod17=9=3^2,\quad
        (T_{139}-c)\bmod139=7=29^2.
\tag{A.134.15}
\]
They are units and local squares.
At $r$ the first entry has the residue of
$h-c=867/382$, whose squareclass is $1146$.
The exact check $(1146/r)=1$ again makes it a
local square.
Thus all these invariants are zero.
The real character is split because $d>0$.
The independent certificate verifies every
valuation and residue in these calculations.

At every prime outside $S$ the coefficients and
the differences $c-i$ are units, and the character
is unramified.
For negative $v_p(T)$ the old equations force
$4\mid v_p(T)=v_p(T-c)$.
For integral $T$ with $T-c$ a unit its invariant
is zero directly.
If $T\equiv c\pmod p$, the old roots give an
actual root tuple in the boundary fibre.
Its Frobenius fixes that tuple and hence fixes
the common character (A.134.8).
These three cases prove zero evaluation outside
$S$ at every old local point, as in A.85.
They also give the finite support of the computation
without an untested prime exception.

The individual characters of $K$ have globally
constant cyclic classes by (A.128.27).
For the singleton character $299$, the norm identity
of A.96 gives
$(T-c,299)=(a_0,299)$.
Their global invariant sums vanish.
Equations (A.134.14)--(15) give sum $1/2$ for
the remaining independent character $d$.
This proves exactly the action of $\gamma_{C,c}$
claimed in the statement.
Global reciprocity
\cite[Theorem 12.1.8]{AgentCTSkoro2019}
excludes a rational point on $\overline C$.
It also excludes any everywhere locally soluble
augmentation $T-c=bY^2$: at its local points
$\alpha$ would equal the constant class $(b,d)$,
whose global invariant sum must be zero.

\emph{The prime repair and its exact splitting-field map.}
The numerical construction has a verified
arithmetic explanation.
Start with the smaller middle coefficient
$a_h^{(0)}=(169/191)139^2$.
Its support places already have the same nonzero
sum, but that finite computation alone does not
prove local solubility everywhere.
For a fresh prime $r$, the desired boundary
root field at $T=h$ for the other two coordinates is
\[
 L=\Q\left(i,\alpha,\beta\right),\qquad
          \alpha^4=191/529,\quad \beta^4=-191/147.
\tag{A.134.16}
\]
Over $\Q(i)$, the valuations at the inert primes
$191$ and $3$ of these two radicands have matrix
$\bigl(\begin{smallmatrix}1&1\\0&-1\end{smallmatrix}\bigr)$.
Its determinant is $-1$.
Kummer independence therefore gives
$N=\operatorname{Gal}(L/\Q(i))=(\Z/4)^2$.
Actual complex conjugation $s$ has order two
and sends every element of $N$ to its inverse.
To check the latter on either root, write
$n(\alpha)=i^a\alpha$ and $s(\alpha)=i^u\alpha$;
then $sns(\alpha)=i^{-a}\alpha$, and the same
calculation applies to $\beta$.
Consequently
\[
 \operatorname{Gal}(L/\Q)=N\rtimes\langle s\rangle,\qquad
 [\operatorname{Gal}(L/\Q),\operatorname{Gal}(L/\Q)]=N^2,
 \qquad L^{\rm ab}=\Q(i,\sqrt3,\sqrt{191}).
\tag{A.134.17}
\]
This is the splitting field's actual abelianization,
distinct from and mapped independently of a
common field of individual fibre factors.
It ramifies only within $\{2,3,7,23,191\}$.

The chosen prime $r$ satisfies $r\equiv1\pmod8$,
is a square modulo every odd prime of
$S\setminus\{r\}$, and splits both quartics in
(A.134.16), as verified by their fourth-power
residue tests.
Thus the smaller support points lift unchanged
after multiplying its middle coefficient by $r^2$.
At $r$ use the final row in (A.134.11).
The two remaining equations lift the nonzero
boundary roots, with unit derivatives.
The small good primes requiring the square class
of $qr$ in this search were
\[
                             5,\ 29,\ 41,\ 233.
\tag{A.134.18}
\]
At each one $(qr/p)=1$.
At each of the other $295$ good primes
$p\le4500$ with $p\equiv1\pmod4$, the retained
data contain a literal centre for the nonsquare
unit twist.  The square unit twist has centre $R$.
For $p\equiv3\pmod4$, the square of any unit
is a fourth power, so $R$ suffices.
This explains the full finite witness construction;
the independent certificate verifies the actual
ratios for the chosen $r$ again.

There are in fact infinitely many prime repairs.
Prescribe $r\equiv1\pmod8$ and $r\equiv1\pmod p$
at every odd smaller support prime, and prescribe
$r\equiv139\pmod p$ for (A.134.18).
These are compatible unit CRT conditions.
Let $\sigma$ be their cyclotomic automorphism
on the field of the resulting modulus.
It fixes $L^{\rm ab}$ in (A.134.17), hence fixes
the intersection with $L$.
The restriction map from the compositum Galois
group to the product of its two factor groups
identifies its image with pairs agreeing on that
intersection.
Thus the pair $(1,\sigma)$ is an actual
automorphism of the compositum.
Chebotarev
\cite[Theorem 8.31]{AgentMilneANT}
gives infinitely many primes with that restriction.
The same $295$ prototypes and the bound
(A.134.13) prove all-place solubility for them.
Their new-prime invariant vanishes because
$h-c$ has squareclass $1146$ supported at the
places where $r$ is a local square.
Explicitly, quadratic reciprocity gives
\[
       \sum_v\operatorname{inv}_v(r,1146)=0,\qquad
       \operatorname{inv}_v(r,1146)=0\ (v\ne r),
             \qquad \operatorname{inv}_r(1146,r)=0.
\]
At $2,3,191$ the first entry is a local square;
at the real place it is positive; at every other
finite place except $r$ both entries are units
and the character is unramified.
These facts prove the middle equality and hence
the last one.
The dyadic sum remains $1/2$.
The explicit prime in (A.134.1) supplies the
executed numerical example used in what follows.

\emph{The fixed boundary fields for completed old coordinates.}
The three boundary quadratic fields are
\[
 \begin{array}{c|c|c}
 i&\text{squareclasses of }(i-j)/a_j,\ j\ne i&B_i\\ \hline
 0&-1146,\ -382&\Q(\sqrt{-1146},\sqrt{-382})\\
 1&382,\ 382&\Q(\sqrt{382})\\
 h&191,\ -573&\Q(\sqrt{191},\sqrt{-573}).
 \end{array}
\tag{A.134.19}
\]
None of these spans contains $-1$.
For the first row the nontrivial product is $3$;
the third has product $-3$; the second is a
positive quadratic field.
The odd valuation at $191$ rules out $-1$ in
the remaining listed classes.
At infinity the two coefficient-ratio
squareclasses are $2,3$, also excluding $-1$.
Multiplying any coefficient by its current
exponent power, with every exponent divisible
by four, preserves all these quadratic fields.
The mixed-exponent root-stabilizer proof
(A.90.3)--(A.90.10) therefore identifies the
common field at each old boundary with exactly
$B_i$ at every subsequent stage.

The initial root-symbol condition needs a sign
calculation; it is not inferred from this
squareclass computation.
At two, the three fourth-power ratios in
(A.134.11) have normalized unit $17$ modulo $32$.
Their actual fourth roots can be chosen with
\[
                  X_0\equiv5,\qquad X_1\equiv3,\qquad
                  X_h\equiv5\pmod8.
\tag{A.134.20}
\]
The fourth-root branch in $1+4\Z_2$ is $5$
modulo eight for such a target; negating it
gives the middle choice.
The resulting invariant table, in the generator
order of (A.134.19), is
\[
 \begin{array}{c|ccc}
    &X_0\text{ against }B_0&X_1\text{ against }B_1&
                                      X_h\text{ against }B_h\\ \hline
 p=2&(1/2,1/2)&0&(0,0)\\
 p=\infty,\quad X_0<0,\ X_1,X_h>0&(1/2,1/2)&0&(0,0).
 \end{array}
\tag{A.134.21}
\]
Here the real point over $R$ is
$(-1,1,1/\sqrt Q)$ in root coordinates.
For a direct dyadic check, a unit $5$ modulo eight
has the unramified quadratic character of $5$.
Its norm group has even valuations.
Both generators $-1146,-382$ have valuation one,
whereas the two generators of $B_h$ are units.
For the middle root, $382$ has squareclass $-2$,
and the identity $3=1^2+2\cdot1^2$ is a literal
norm proving its zero entry.
These prove every dyadic value in (A.134.21).
The real values follow from the displayed signs.

At $3$ and $191$ use centre $R$.
All three roots are units.
Since these primes are $3$ modulo four, exactly
one of the two signs of a fourth root has square
unit residue; choose that sign.
Its unit is then a local square, so every
boundary-field symbol is zero.
At all other witnessed primes the $B_i$ are
unramified.
The certificate checks the exact conditions
\[
       v_p(X_i)\,\chi(\operatorname{Frob}_p)=0
             \quad\text{for each generator of }B_i.
\tag{A.134.22}
\]
There are $3040$ such checks.
At larger good primes choose the unit-root
points supplied by (A.134.13); all these symbols
are zero because the roots are units.
Thus the only nonzero initial root symbols are
the two dyadic entries and their matching real
entries in (A.134.21).
Their \emph{global sums} are zero.
The two other root coordinates already have
zero symbols at each place individually.

\emph{One signed refinement establishes the full invariant.}
We now give the required initial modification
of the exact A.90 construction.
Start with exponents $(4,4,4)$ and refine
coordinate zero to exponent eight.
Take a finite controlled set $S_0$ containing
all the above bad places, all primes up to the
new mixed-cover bound (A.90.19), and all the
ramified primes of the boundary splitting field.
At the additional controlled primes take
unit-root old points by (A.134.13).
Their root symbols vanish.
All old points avoid $T=c$.

Write $x_p=X_0(P_p)$ at these places and use
the negative real root $x_\infty=-1$.
Set
\[
                 d_0=-\prod_{p\in S_0}p^{v_p(x_p)}<0,
                  \qquad v_p^*=x_p/d_0.
\tag{A.134.23}
\]
Choose precisions $M_p>2v_p(2)+4$, large enough
that $\mathcal R=\Q(\mu_M)$ contains $B_0$,
where $M=\prod_{p\in S_0}p^{M_p}$.
Let $\sigma\in(\Z/M)^\times$ be the CRT class
$\sigma\equiv v_p^*\pmod{p^{M_p}}$.
For every character of $B_0$, its total
evaluation on the $x_p$ and $x_\infty$ is zero
by (A.134.21)--(22).
Subtract the reciprocity sum of the rational
number $d_0$.
Outside $S_0$, $d_0$ is a unit in an unramified
character field, and $x_\infty/d_0>0$.
It follows that
\[
                 \prod_{p\in S_0}\operatorname{rec}_p(v_p^*)|_{B_0}=1,
                         \qquad \sigma|_{B_0}=1.
\tag{A.134.24}
\]
The second equality uses the inverse-unit
cyclotomic reciprocity convention, exactly as
in (A.90.21)--(23).
This is where the negative real sign cancels
the dyadic symbols.

The boundary algebra at $X_0=0$ has radicands
$(0-j)/a_j$ for $j=1,h$.
Its common field is $B_0$ by (A.134.19) and
the mixed-exponent calculation.
The exact stabilizer-image equality (A.90.25)
expresses $\sigma$ as a product of images of
actual boundary root-tuple stabilizers.
Chebotarev supplies distinct fresh primes
$q_l$ realizing those images, and an additional
prime $q_*$ splitting the whole compositum.
Thus
\[
 \begin{gathered}
       u=d_0q_*\prod_lq_l<0,\qquad
             u/x_p\in1+p^{M_p}\Z_p\quad(p\in S_0),\\
       X_0=uY_0^2,\qquad X_1=Y_1,\qquad X_h=Y_h,\qquad
       T=T,\\
       a'_0=a_0u^4,\quad n'_0=8,\qquad
       (a'_1,n'_1)=(a_1,4),\quad
       (a'_h,n'_h)=(a_h,4).
 \end{gathered}
\tag{A.134.25}
\]
This is an actual morphism of old covers by substitution.
At every controlled finite prime, Hensel's lemma
gives $Y_0=(x_p/u)^{1/2}\in1+p^3\Z_p$.
It is itself a local square.
At the real place $x_\infty/u>0$, and choose the
positive root.
All symbols of the new coordinate are now
individually zero; the other roots were retained
and have their previous individual zero symbols.

At a fresh coefficient prime, set $Y_0=1$ and
$T=a_0u^4$.
The fixed boundary tuple supplies nonzero
residue roots for the other two equations,
which lift because the prime divides no exponent.
All new roots are units and all fixed $B_i$
are unramified there, so their symbols are zero.
The point avoids $c$ and the other old punctures:
its distance from zero has positive valuation,
whereas the relevant boundary differences are
units.
Outside the enlarged controlled set, the
mixed-cover bound gives unit-root points
avoiding also the fibre $T=c$, exactly as in
(A.90.18)--(19).
Thus (A.134.25) establishes the full
pointwise symbol-zero invariant required for
all subsequent refinements.

\emph{Cofinal refinement and the actual completed tuple.}
Continue by refining the other two coordinates
by the factor two, then refine all three
coordinates in cycles with factors
$3,4,5,\ldots$.
At a stage, for the chosen coordinate $i$ and
factor $m$, the literal map is
\[
 \begin{gathered}
 X_i=uY_i^m,\quad X_j=Y_j\ (j\ne i),\quad T=T,\\
                    a_i'=a_i u^{n_i},\qquad n_i'=mn_i.
 \end{gathered}
\tag{A.134.26}
\]
Use the positive real roots after (A.134.25).
The individual zero symbols, fixed quadratic
fields (A.134.19), and actual boundary common
fields give the precise CRT condition of
(A.90.20)--(23).
Its stabilizer product (A.90.24)--(26) constructs
a positive rational $u$, including one extra
splitting prime of valuation one.
At controlled places the root $Y_i$ is principal
and square, while $T$ and all other roots are
retained exactly.
At fresh primes take $Y_i=1$ and
$T=i+a_i u^{n_i}$; the boundary tuple and unit
derivatives give the other roots.
At remaining good primes the bound
(A.90.19) gives unit-root points avoiding $c$.
These are precisely the three local cases
proved in (A.90.27)--(28).
All their hypotheses have been verified above
for the present fields and signed initial step.
Enlarge each controlled set to contain the
primes at most the stage index.
Each prime's $T$ is therefore fixed permanently
after a finite stage.

After the cycle with factor $m$ the common
exponent is $4m!$.
The coefficient power identity in (A.134.26)
preserves its preceding class, so these
rational coefficients determine compatible
classes $A_i$ modulo every positive integer.
At every fixed finite prime and every modulus
$N$, a sufficiently late retained equation has
exponent divisible by $N$ and reads
$t_p-i=a_{i,j}X_{i,j,p}^{n_{i,j}}$.
It gives the exact localization equalities
(A.134.5) on taking all moduli.
No convergence assertion for changing local
centres is needed: their centres have been
retained literally.
The initial dyadic point gives $t_2=31/2$.
At the real place every coefficient has
remained positive, and $T=R$ was retained.

The extra splitting prime at every refinement
has coefficient valuation $n_i>0$ in that
coordinate and zero in the others.
It is immediately controlled with unit roots.
Every later coefficient multiplier has valuation
zero there.
These exact positive integer valuations persist
in the completed class.
Infinitely many distinct such primes are
introduced for each coordinate, proving the
infinite-support assertion.

Finally the maps (A.134.25)--(26) make every
retained $t_p$ an actual point of the initial
quartic $C$.
The invariant of $(t_p-c,d)$ is therefore $1/2$
at two and zero at every other finite place,
by the already proved local constancy.
The real character is split because $d>0$.
If a rational squareclass $b$ matched
$[t_p-c]$ at all finite places, the constant
quaternion $(b,d)$ would have exactly these
invariants.
Their sum contradicts global reciprocity.
This excludes the squareclass and hence every
completed new coordinate, completing the proof.

\emph{An effective convention specifying one completed tuple.}
Use the retained rational centres at the $611$
listed finite primes.
At two take exactly the fourth-root branches
in (A.134.20); at $3$ and $191$ take the unique
fourth-root branch with square unit residue.
At every other listed prime use the least
nonzero residue root in each coordinate after
removing its indicated uniformizer power,
and the unique simple-root Hensel branch.
The unramified root-symbol calculation is
independent of this last choice.
Retain the real roots in (A.134.21).

Order rational primes increasingly and all
finite residue tuples and root-label groups
lexicographically.
At an additional controlled good prime choose
the first ordinary residue point of the
old curve, take its least nonnegative integral
centre, and take its indicated Hensel branches.
If necessary, increase that centre by the
least nonnegative multiple of the prime
which avoids $c$.
At an outside good prime use the first
residue point on the open curve with both
old boundary and $T=c$ removed.
Choose the least integer precisions satisfying
the stated Hensel and conductor conditions.
At each refinement compute the actual finite
mixed root stabilizers by (A.90.5)--(9),
enumerate finite words in their restriction
images, and take the first word with the
prescribed CRT product.
For its successive Frobenius classes take
the least admissible distinct primes, and
the least additional completely split prime
for the required extra factor.
The initial sign is the negative one in
(A.134.23); all later real roots and factors
use the positive choices proved above.
The close-to-one Hensel root at a controlled
place is the specified principal branch.

The point-count, exact stabilizer-generation,
and Chebotarev arguments prove termination
of each finite selection.
Together with the factorial cycle these
conventions specify one actual completed
tuple.  The retained computations certify
the finite starting data; no execution of
the infinite prime recursion is asserted.
\end{proof}

The numerical search and its exact data are retained in
\path{code/harmonic_twist_repair_search.py} and
\path{code/harmonic_twist_repair_data.json}.
The independent standard-library certificate
\path{code/harmonic_twist_repair_certificate.py}
checks all $611$ finite centres, $1833$ fourth-power
conditions, the complete support and invariant table,
the degree-sixteen common field, the $3040$ unramified
initial root-symbol conditions, the signed dyadic
and real cancellation, the splitting-field group,
the $295$ nonsquare prototypes, and the good-prime
bound.
Its full output is retained.
The separate
\path{code/harmonic_root_symbol_independent_audit.py}
independently checks the initial root symbols.
